\hspace{3mm}

\centerline{\bf \huge 9 класс}
\centerline{\normalsize{\it Работа рассчитана на \bf{180} минут}}

\hspace{3mm}

\problem{9.1} Два друга записали на доске следующие суммы: Адьян $2023^{2023} + 2021^{2021}$, а Церен $2023^{2021} + 2021^{2023}.$ Определите, у кого из друзей получилась большая сумма?\

\hspace{3mm}

\problem{9.2} На городской этап олимпиады по математике пришли 2023 ученика из разных школ Элисты, при этом у любых двух участников олимпиады разное количество знакомых среди других участников. Найдётся ли ученик, который будет знаком со всеми другими участниками городской олимпиады по математике? 

\hspace{3mm}

\problem{9.3} Учитель выписал на доску все натуральные числа от 1 до 2023 и предложил Баиру каждую минуту выполнять следующую операцию до тех пор, пока у него не останется одно число на доске: стирать любые два числа $a$ и $b$ с доски и записывать вместо них одно число, равное $2022\cdot(ab)^{2023} + 1^{2023}$. Баир согласился и быстро справился с заданием. Какой остаток при делении на два даёт последнее оставшееся число на доске?

\hspace{3mm}

\problem{9.4} В параллелограмме $ABCD$ проведены диагонали $AC$ и $BD$, которые пересекаются в точке $O$. Точка $M$ - середина стороны $CD$. Точка $K$ делит диагональ $AC$ в отношении $AK$:$KC$ = 1:2. Докажите, что прямые $MO$ и $DK$ пересекаются на середине стороны $AB$.

\hspace{3mm}

\problem{9.5 } Дан правильный 97-угольник. У Алтаны есть 10 банок с красками различных цветов. Она посчитала, сколькими способами может покрасить вершины этого 97-угольника своими красками, но те раскраски, которые переходят друг в друга поворотами, Алтана считает одинаковыми. Какое число у нее получилось? 
Указание: 97 - простое число
